\chapter{Experimental Results}
\label{ch:results}

% --------------------------------------------------------------------------
\section{Supervised Fine-Tuning}
\label{sec:results-sft}
% --------------------------------------------------------------------------

\subsection{Training Dynamics}
\label{sec:sft-training}

Qwen2.5-Coder-1.5B~\cite{qwen25coder} was fine-tuned on the 27,514-example SFT
dataset using rank-16 QLoRA on the Q, K, V, and O projections of all attention
layers.
Training ran for 3 complete epochs (9,288 gradient steps) with an effective batch
size of~8 via gradient accumulation, maximum sequence length of 768~tokens,
learning rate $2 \times 10^{-4}$ with AdamW 8-bit, and BF16 precision with
FlashAttention-2.

The evaluation loss decreased from 0.637 at step~0 to 0.557 at step~9,288.
The loss curve stabilises from the end of epoch~2 (approximately step~6,000)
onward, with train and eval loss tracking closely throughout: no significant
overfitting is observed.
The 3-epoch checkpoint is selected as the SFT model for all downstream
evaluation and RL training.

\subsection{Pass-Rate Evaluation}
\label{sec:passrate}

The fine-tuned model is evaluated against the zero-shot Qwen2.5-Coder-1.5B base
model.
Both models generate 100 programs from the same prompt template at temperature
0.9 with greedy-decoding fallback for programs that fail to parse.
Each generated program is passed through the pure-Python BPF encoder; programs
that encode successfully are submitted to the \texttt{kcov\_validator} binary
inside the KCOV VM; the verifier verdict is recorded.
The pass-rate is defined as the fraction of the 100 generated programs that
receive an ACCEPTED verdict from the Linux~6.8 verifier.

\begin{table}[h]
\centering
\caption{SFT pass-rate: fine-tuned model vs.\ zero-shot baseline}
\label{tab:passrate}
\begin{tabular}{lrrrr}
\toprule
Model                    & $N$ & Encoded & Accepted & Pass-rate \\
\midrule
curated-merged (SFT)     & 100 & 73       & 60       & 60.0\%    \\
zero-shot (base model)   & 100 &  1       &  1       &  1.0\%    \\
\bottomrule
\end{tabular}
\end{table}

The fine-tuned model achieves a verifier pass-rate of 60.0\%, compared to 1.0\%
for the zero-shot base model: a 60-fold improvement (Table~\ref{tab:passrate}).
These results are sourced from \texttt{results/passrate\_summary.csv},
\texttt{results/passrate\_run\_curated.log}, and
\texttt{results/passrate\_run\_zeroshot.log}.

The 13 programs that encoded successfully but were rejected (73 encoded, 60
accepted) contain syntactically valid BPF bytecode that the verifier rejects for
semantic reasons: out-of-bounds memory access, invalid type coercions, or
constraint violations that the model generates plausibly but cannot fully satisfy.
This gap is expected: the SFT objective is next-token prediction over the training
distribution, not verifier satisfaction; the model acquired BPF syntax and
register-naming conventions, but not the complete semantics of the verifier's
type and bounds analysis.
These partially-correct programs --- valid enough to encode, but rejected at the
semantic level --- are in fact the most interesting inputs for RL training: they
reach the verifier and trigger its analysis logic, providing coverage signal even
in the absence of acceptance.

An extended evaluation over 400 generated programs
(\texttt{results/passrate\_curated-merged-400.csv}) yields a pass-rate of 54.8\%
(219 accepted of 400 generated, 226 encoded), confirming that the 60\% figure
from the primary evaluation is representative and not inflated by sample variance.

These results answer RQ1 affirmatively: a 1.5-billion-parameter code language
model can acquire sufficient knowledge of BPF verifier syntax to pass the
verifier 60\% of the time, using parameter-efficient fine-tuning on consumer
hardware (8~GB VRAM).

% --------------------------------------------------------------------------
\section{RL Training: Run 1 (\texttt{rl\_grpo\_v2})}
\label{sec:results-rl}
% --------------------------------------------------------------------------

\subsection{Coverage Exploration}
\label{sec:coverage}

RL run~1 was executed with GRPO configuration: $G = 4$ completions per prompt,
KL penalty $\beta = 0.01$, maximum completion length 600 tokens, running on the
local RTX~4070~Laptop GPU (8.59~GB VRAM).
The run executed for 8,370 gradient steps.

Two coverage metrics are reported, and they must be interpreted carefully as they
measure distinct quantities.

\emph{Total unique PCs} (1,638): the size of the cumulative coverage frontier
at the end of run~1, counting every distinct kernel PC address that appeared in
any KCOV trace across all evaluated programs
(source: \texttt{results/rl\_pc\_set.json}).
This set includes PCs reached during the SFT pass-rate evaluation (which
established a baseline of approximately 1,501 unique PCs before RL training
began) as well as PCs discovered during the RL run itself.

\emph{Incrementally new PCs} (137): the number of kernel PC addresses that were
seen for the first time during RL training, i.e., PCs not present in the
coverage frontier at the start of the run
(source: \texttt{results/rl\_analysis\_pcs.csv}).
This is the direct measure of how much additional exploration the RL phase
contributed.
The README records 138 new PCs, which is an off-by-one artefact from the
boundary-counting convention; the authoritative figure from the CSV is~137.

Coverage growth was highly concentrated.
No new PCs were found for the first 97 batches.
The first discovery event occurred at batch~97 (4 new PCs);
by batch~1,241, the cumulative count had reached~135 new PCs through a series
of bursts.
After batch~1,241, only two further PCs were ever found (at batches~3,036
and~4,873), and no new coverage was discovered in the remaining $\approx$6,500
batches of the run.

\subsection{Plateau Diagnosis}
\label{sec:plateau-results}

Analysis of \texttt{results/rl\_analysis\_tiers.csv} reveals the reward collapse
in precise terms.
Across 7,823 reward-evaluation batches in segment~1, the within-group reward
standard deviation $\sigma_r$ was non-zero in only 41 batches (0.5\%).
For the remaining 7,782 batches (99.5\%), $\sigma_r = 0$.
The 41 non-zero batches correspond almost exactly to the 9 discovery events
listed above: a batch that finds new PCs assigns reward~1.0 to the completing
program and 0.1 to the others, producing $\sigma_r > 0$ and a non-zero gradient.
Once the discovery events cease after batch~1,241, no further non-zero $\sigma_r$
occurs except for rare isolated batches ($\sigma_r = 0.05$) that reflect a single
completion with a marginally different reward in an otherwise uniform group.

The GRPO gradient is proportional to the advantage $A_i = (r_i - \bar{r}) / \sigma_r$
(Equation~\ref{eq:grpo-advantage}).
When $\sigma_r = 0$, the gradient is identically zero: the model parameters are
not updated, and the training run degenerates into a consume-compute loop with no
learning.
This is the gradient starvation failure mode documented in~\cite{grpo-starvation}.

The verdict breakdown illuminates the structural cause.
Over the entire run, 31,125 of 31,258 completions (99.7\%) were accepted by the
verifier; only 102 (0.3\%) were rejected, and 31 (0.1\%) failed to encode.
The near-total absence of rejected programs is a direct consequence of SFT: the
model was trained on a prompt with a \texttt{Status:~VALID} header, and it
learned to generate programs that satisfy the verifier almost all of the time.
During RL inference, this behaviour persists.
As a result, almost every GRPO group consists of four accepted programs, all of
which receive reward~0.1 once the coverage frontier is saturated --- producing
uniform reward and zero gradient.

A secondary factor is the KCOV bug fixed in commit~\texttt{c8a9d02}: even the
small minority of rejected programs would have appeared identical to the reward
function (all returning \texttt{pcs:~[]}) prior to the fix, further suppressing
variance in any group that contained a rejected completion.
Because both the prompt bias and the KCOV bug coexisted in run~1, their individual
contributions to the gradient collapse cannot be separated from the run~1 data
alone.

% --------------------------------------------------------------------------
\section{Reward Redesign Validation}
\label{sec:reward-validation}
% --------------------------------------------------------------------------

Following the bug fix in commit~\texttt{c8a9d02}, the corrected
\texttt{kcov\_validator} binary was exercised in a 100-step comparison run
(Config~A: $G=4$, max\_completion\_length=600, $\beta=0.01$, neutral prompt
without a \texttt{Status:~VALID} header).
Source: \texttt{results/comparison\_summary.txt} and
\texttt{results/grpo\_completions.log}.

The comparison run generated 400 completions over 100 steps, with the following
verdict breakdown: 46 accepted (11.5\%), 253 rejected (63.2\%), and 101 encode
failures (25.2\%).
The contrast with run~1's verdict distribution (99.7\% accepted) is striking and
confirms the prompt-bias hypothesis directly: removing the \texttt{Status:~VALID}
header shifts the model's output from nearly all-accepted to predominantly
rejected, demonstrating that the SFT prompt was indeed anchoring the model toward
valid-program generation during RL inference.

The 253 rejected programs in the comparison run returned non-empty KCOV traces
from the fixed binary, as evidenced by the reward function assigning
differentiated depth scores to those programs.
This validates the structural fix: after commit~\texttt{c8a9d02}, rejected
programs are treated identically to accepted programs for the purpose of coverage
attribution.
The comparison run found 2 new kernel PCs (from a 1,501-PC baseline), indicating
that the fixed pipeline is capable of coverage-guided exploration even in a
predominantly-rejected regime.

The depth-based verdict-blind reward (Equations~\ref{eq:depth}--\ref{eq:reward-v2})
is designed to maintain non-zero within-group reward variance in this regime.
With rejected programs now returning varied PC trace lengths --- rather than
uniform empty traces --- the depth component assigns different scores to each
completion based on how deeply the verifier processed that program, providing a
gradient signal even when all $G$ completions are rejected.
Full validation of this redesign under sustained training requires run~2.

% --------------------------------------------------------------------------
\section{Run 2: Planned Experiment}
\label{sec:run2}
% --------------------------------------------------------------------------

Run~2 is the next experimental step, planned but not executed within the scope
of this thesis.
Its configuration is:

\begin{itemize}
  \item \textbf{Reward function:} depth-based verdict-blind reward
        (Section~\ref{sec:reward-redesign}, \texttt{ml/reward.py}).
  \item \textbf{Prompt:} neutral header with no \texttt{Status:~VALID} prefix,
        removing the valid-program attractor identified in run~1.
  \item \textbf{Group size:} $G = 8$ completions per prompt, doubling the
        within-group sample size relative to run~1 for a more reliable
        $\sigma_r$ estimate.
  \item \textbf{Max completion length:} 1,200 tokens, allowing the model to
        generate longer programs that may exercise deeper verifier paths.
  \item \textbf{Compute:} Colab Pro (A100 GPU) via the remote reward pipeline
        (Section~\ref{sec:remote-reward}).
\end{itemize}

The remote reward pipeline is fully implemented and locally validated: the
\texttt{tools/reward\_server.py} FastAPI server, the Cloudflare Quick Tunnel,
and the \texttt{ml/train\_grpo\_colab.ipynb} notebook are all in a
run-ready state.
The experimental phase of this thesis closed before the graduation deadline
without executing run~2.

The expected outcome is that the depth-based reward maintains $\sigma_r^2 > 0$
throughout training, avoiding the gradient starvation that halted run~1, and
that the neutral prompt widens exploration by allowing the model to generate
programs across the full range of verifier outcomes rather than converging on
the valid-program mode.
Whether these design changes are sufficient to drive sustained coverage growth
beyond the 137 new PCs found in run~1 is the open research question that
run~2 will address.
